El 29 de novembre de 2018, el nanosatèl·lit \textit{CubeCat-1}, desenvolupat per estudiants i investigadors de la Universitat Politècnica de Catalunya (UPC), es va llançar a l'espai des de la base espacial de Sriharikota, a la costa est de l'Índia, dins d'un coet de l'agència espacial índia ISRO.

El \textit{CubeCat-1} té una massa d'1,30~kg i orbita a 530~km de la superfície de la Terra.

\figura[0.4]{fig1.png}
{\centering \textsc{Font}: \textit{https://www.upc.edu}.\par}

\begin{apartat}{1}
Calculeu el període orbital del \textit{CubeCat-1} i indiqueu el nombre de voltes completes que fa cada dia al voltant de la Terra.
\end{apartat}

\begin{apartat}{1}
Quin és el pes del nanosatèl·lit en la seva òrbita?
\end{apartat}

\dades{%
  $G = 6,67\times 10^{-11}\ \mathrm{N\,m^2\,kg^{-2}}$.\\
  $M_\mathrm{Terra} = 5,98\times 10^{24}\ \mathrm{kg}$.\\
  $R_\mathrm{Terra} = 6,37\times 10^{6}\ \mathrm{m}$.}
